\section{A genus budget and a small attachment network}
\label{sec:attachment-kernel}

The initial torus-boundary reduction is the smallest instance of a more
general structural fact.  Its useful parameter is the topology of the
boundary, rather than the number of zero-handles in a presentation.  We
give the argument in a form that separates a topological theorem from the
combinatorial certificate implemented in this package.

\subsection{Cut pieces and the unused-handle graph}

Let $M$ be a connected compact orientable irreducible three-manifold with
nonempty boundary.  Suppose a finite handle presentation has been given
and a set $U$ of one-handles is disjoint from the attaching regions of all
higher-index handles.  Their cocores form disjoint properly embedded
discs.  Cut along all these discs.  Let $K_v$, for $v\in V(Q)$, be the
connected pieces, and put one edge of $Q$ for every original handle in
$U$, with endpoints its two incident pieces.  Loops and parallel edges
are retained.  The graph $Q$ is connected, and
\[
 q=|V(Q)|,\qquad e=|E(Q)|,\qquad \beta=e-q+1.
\]
The index $\beta$ is a graph cycle rank, not the first Betti number of the
whole manifold.  Write
\[
 \delta(M)=1-\chi(M),\qquad d_v=1-\chi(K_v).
\]
When $\partial M$ is connected of genus $g$, orientability and
Poincar\'e--Lefschetz duality give $\delta(M)=g$.  For boundary components
of genera $g_1,\ldots,g_b$, the corresponding formula is
$\delta(M)=1-b+\sum_i g_i$.

\begin{lemma}[Irreducibility and Euler labels after disc cutting]
\label{lem:disc-euler}
Every $K_v$ is irreducible and has nonempty boundary.  It is a ball if and
only if $\chi(K_v)=1$.  Every other piece satisfies $\chi(K_v)\leq0$.
\end{lemma}
\begin{proof}
An interior sphere in a cut piece is also an interior sphere in $M$, so
it bounds a ball in $M$.  The cutting discs are disjoint from the sphere.
Each disc is connected and has boundary on $\partial M$, outside that
interior ball.  It therefore lies outside the ball.  Consequently the
ball survives inside the same cut piece.  Disc cutting cannot create a
closed component, which proves the first assertions.

A connected irreducible three-manifold with a sphere boundary component
is a ball: push that component into the interior and apply irreducibility;
the side away from the original boundary is a ball, and restoring the
spherical collar gives a ball.  Thus a nonball $K_v$ has only boundary
components of genus at least one.  The identity
$2\chi(K_v)=\chi(\partial K_v)$ now implies $\chi(K_v)\leq0$.
Conversely a ball has Euler characteristic one, and the two possibilities
exhaust the pieces.
\end{proof}

This lemma uses actual irreducibility of the source.  An abstract integer
label equal to one, a homology computation, or a sphere boundary alone
does not establish the conclusion.  The executable reducible controls in
Section~\ref{sec:evidence} show why the distinction matters.

\begin{theorem}[Exact unused-handle budget]
\label{thm:genus-budget}
Under the preceding hypotheses,
\begin{equation}
 \boxed{\quad
 \delta(M)=\beta+\sum_{v\in V(Q)}d_v
          =\beta+\sum_{K_v\text{ nonball}}\bigl(1-\chi(K_v)\bigr).
 \quad}
 \label{eq:genus-budget}
\end{equation}
In particular, if $s$ is the number of nonball pieces, then
$s+\beta\leq\delta(M)$.  For connected boundary this is
$s+\beta\leq g$.
\end{theorem}
\begin{proof}
Recover $M$ by attaching the $e$ one-handles to the disjoint union of the
$K_v$.  Each such attachment decreases Euler characteristic by one, so
$\chi(M)=\sum_v\chi(K_v)-e$.  Substituting $e=\beta+q-1$ gives
\[
 1-\chi(M)=\beta+\sum_v(1-\chi(K_v)).
\]
By Lemma~\ref{lem:disc-euler}, a ball contributes zero and every nonball
contributes an integer at least one.  This proves the equality and bound.
\end{proof}

The identity does more than count components: it distinguishes the two
possible sources of boundary complexity, cycles in the attachment graph
and topological complexity inside the retained pieces.  Neither term can
be large while the boundary genus remains small.

\subsection{Absorbing ball pieces along a rooted forest}

For the following topological replacement, a network means the graph
together with its vertex manifolds and its labelled attaching discs.
The graph without these interfaces is not a complete manifold encoding.
No boundary pattern or prescribed embedded curve is part of the replacement
contract.

\begin{lemma}[Absorption of a ball vertex]
\label{lem:absorb-ball}
Suppose a ball vertex $B$ is joined to a different vertex manifold $K$ by
a chosen one-handle.  The union of $K$, this handle, and $B$ is
homeomorphic to $K$.  Under such a homeomorphism all other attaching discs
on $B$ can be moved to mutually disjoint subdiscs of a neighbourhood of
the chosen attaching disc on $\partial K$.  The homeomorphism can be taken
to fix the complement of a chosen three-dimensional boundary collar in $K$.
\end{lemma}
\begin{proof}
A ball attached to one end of a one-handle is a ball with a distinguished
disc at the other end.  Gluing this ball to $K$ along that disc adds a
boundary collar.  Removing the collar gives the asserted homeomorphism.
The remaining attaching discs are a finite labelled collection of
disjoint discs on the boundary of the added ball, away from its
distinguished disc.  A homeomorphism of the boundary sphere carries
them into any prescribed disjoint subdiscs in the corresponding collar
patch.  It extends over the ball.  This argument preserves labels and
orientation, including when some of the other handles have both ends on
$B$.  It does not preserve an arbitrary additional boundary pattern.
\end{proof}

\begin{theorem}[Genus-controlled attachment kernel]
\label{thm:attachment-kernel}
Assume the hypotheses of Theorem~\ref{thm:genus-budget}, and suppose the
attaching discs have been supplied.  There is a topologically equivalent
unmarked attachment network with the following sizes.
\begin{enumerate}[label=(\roman*)]
\item If $s\geq1$, the vertex manifolds are exactly the $s$ nonball
pieces, and the number of one-handles is
\begin{equation}
             e'=\beta+s-1\leq\delta(M)-1.
             \label{eq:kernel-edge-count}
\end{equation}
Thus the network has at most $g$ vertex manifolds and $g-1$ handles when
the original boundary is connected of genus $g$.
\item If $s=0$, the network consists of one ball and $\beta=\delta(M)$
one-handles.  In particular $M$ is a handlebody of genus $\delta(M)$.
\end{enumerate}
The source-to-kernel combinatorial correspondence has a rooted-forest
certificate of $O(q+e)$ records, constructible and checkable in polynomial
bit time.  This is a bound for the attachment network; it is not a bound
for the internal presentation size of the retained nonball pieces.
\end{theorem}
\begin{proof}
If $s\geq1$, choose a spanning forest $F$ of $Q$ with exactly one nonball
vertex in each tree.  One constructive method starts with all vertices
separate, and processes the edges, joining two distinct components unless
both already contain a nonball vertex.  At most one root is ever present
in a component.  At termination, any component with no root would have
an edge to another component because $Q$ is connected.  That edge could
not have been rejected by the rule.  Hence every final component has a
root and there are exactly $s$ trees.

Root each tree at its nonball vertex.  Starting at the leaves, absorb
every other vertex and its parent handle using
Lemma~\ref{lem:absorb-ball}.  Every absorbed vertex is a ball.  Incidences
of handles outside $F$ are retained, with their attaching discs carried
into the appropriate parent patch.  Exactly $q-s$ handles are removed,
leaving
$e-(q-s)=\beta+s-1$.  The bound follows from
Theorem~\ref{thm:genus-budget}.

If every piece is a ball, choose one vertex as the root of a spanning
tree and carry out the same procedure.  The remaining ball has
$e-(q-1)=\beta$ one-handles, the defining presentation of an orientable
handlebody of genus $\beta$.  Equation~\eqref{eq:genus-budget} gives
$\beta=\delta(M)$.

For the combinatorial certificate retain the original edge identifiers,
the edges of $F$, the root of every source vertex, and one parent pointer
per nonroot vertex.  Every surviving handle is still an original
nonforest edge.  Its new ends are the roots of its old ends.  A checker
verifies the source hash, all forest edges, acyclicity, unique roots,
coverage of every source vertex, every remaining incidence, and the Euler
and cycle-rank equalities.  Sorting and union--find give
$O((q+e)\log(q+e))$ comparison work, with binary costs charged separately.
The independent checker can traverse the supplied rooted trees directly.
\end{proof}

\subsection{Source addresses and the meaning of a small kernel}

An endpoint originally on an absorbed vertex is rerouted along its parent
chain.  These chains must be shared.  Explicitly writing the complete
route for every endpoint of a long chain can require quadratic space;
one parent pointer per source vertex suffices.  The implementation returns
this forest.  After a one-time $O(q)$ parent-index construction, it permits
expansion of one requested route in time proportional to its length.
The convenience API pays that indexing cost if no index is supplied.
It never expands every route as part of
kernel construction or verification.

There are three distinct outputs here.  The operational graph has the
small size in Theorem~\ref{thm:attachment-kernel}.  Its source provenance
still has $O(q+e)$ records.  An explicit homeomorphism from a triangulated
ball piece to a standard ball is a further geometric output and is not
constructed by the graph routine.  The topological replacement is
justified by the preceding lemmas once the vertex pieces and source
irreducibility are authenticated.  A later patterned-hierarchy consumer
must additionally construct and verify the required disc charts and
pattern transport.  The initial torus compiler avoids this extra task by
retaining its unique core subpresentation directly.

\begin{proposition}[Sharp counts for this replacement contract]
\label{prop:kernel-sharp}
The bounds $s\leq g$ and $e'\leq g-1$ in the nonball case are attained
simultaneously for every $g\geq1$.  In the all-ball case the one-vertex
graph requires exactly $g$ loops.
\end{proposition}
\begin{proof}
Take $g$ solid tori and connect them by $g-1$ one-handles in a tree.  The
result is a handlebody of genus $g$, hence irreducible.  Regard the
connecting handles as $U$.  Its pieces are the $g$ solid tori, each with
Euler characteristic zero, so $s=g$ and $\beta=0$.  Arbitrarily many ball
vertices can be inserted along the tree without changing the example.
All $g$ specified nonball vertex manifolds remain under the stated
replacement contract, and connecting them requires at least $g-1$
handles.  If every piece is a ball, cycle rank is $g$ and a one-vertex
graph has exactly one independent cycle per loop.
\end{proof}

This sharpness is relative to retaining the specified nonball blocks.  A
handlebody could of course be given another global presentation.  More
generally, a knot exterior of arbitrarily high presentation complexity can
be the sole nonball vertex when $g=1$.  A one-vertex attachment graph
therefore gives no uniform bound on the difficulty of recognizing that
vertex manifold.  This is why the result improves a hierarchy primitive
without proving general quasi-polynomial unknot recognition.

\subsection{Implementation and independent graph replay}

\code{attachment\_kernel.py} implements the rooted-forest producer.
\code{attachment\_kernel\_verify.py} imports no producer code and checks
the complete finite certificate.  The adapter
\path{kernel_from_decomposition} converts the decomposition format from
\code{unused\_handles.py}; it does not call that module's geometric
verifier.  The caller must first authenticate the decomposition against
the original source.  In particular, the graph certificate's hash binds
only the vertex Euler labels and edge incidences.  It does not bind a
triangulation of each vertex manifold or its attaching-disc charts.
A topological consumer must retain the separate verified source packet
and, when invoking ball absorption, its irreducibility justification.
Loops, parallel handles, isolated one-vertex networks, zero genus and
nonball deficits larger than one are part of the supported graph model.
The tests include every connected labelled simple graph on at most five
vertices under three root assignments, random multigraphs, sharp families,
and altered proof records.
